\section{Discussion}

\subsection{Main interpretation}

The current evidence supports partner-local affective precision as computational infrastructure for social policy deployment. H1 shows that precision is more closely tied to predictive reliability than realized payoff. H2 shows that decoupling beta from policy precision can preserve similar belief readouts while changing entropy and payoff. H4 shows that partner-choice distributions can move before payoff moves. H3 shows that the same precision channel can become harmful when social volatility is abrupt.

This combination is stronger than a monotonic reward-improvement story. A generic reward booster would predict payoff gains and would not naturally explain why precision follows surprise more strongly than reward. A pure belief-accuracy story would predict improvements in type or stance readouts, but H2 shows that belief readouts can remain similar while deployment changes. A pure caution story would predict that making the agent more cautious repairs betrayal failure, but the H3 sensitivity results show that generic caution can increase entropy while reducing payoff.

The model therefore occupies a middle position. Partner-local affective precision is useful because it compresses model fitness by partner and gates action on the basis of that confidence. It is dangerous because confidence can be misplaced during abrupt model mismatch. Affect is not a guarantee of safer recovery; it is a mechanism that can improve or impair deployment depending on whether the underlying partner model is calibrated to the current regime.

\subsection{Why this is not reward, trust, or deeper planning}

The H1 dissociation is the central argument against reward averaging. In the 30-seed H1 confirmation, affect had lower total payoff than no affect by $7.53$ points, with a confidence interval spanning zero, yet precision tracked surprise more strongly than payoff. This is exactly what a model-fitness account predicts. The affective state says ``my model of this partner is reliable'' rather than ``this partner is good for me.''

The mechanism is also not a literal trust scalar. Trust can mean willingness to rely on a partner, expected benevolence, or learning-rate modulation. Here, beta modulates policy precision. A predictable adversary can have high model precision but low approach value. A future model could add a separate trust-as-learning-rate variable, but the current implementation intentionally keeps trust and affective precision distinct.

The mechanism is not simply deeper planning. The current architecture uses action-dependent stance dynamics so planning depth is meaningful, but the manuscript's claims concern partner-local precision as an orthogonal deployment signal. Policy precision can change how the agent uses a policy set; it is not another step of lookahead. This is why policy openness matters: if the policy posterior is already saturated, precision changes may be invisible. If several policies remain plausible, precision can strongly alter behavior.

\subsection{Boundary conditions as evidence}

The H3 betrayal result should be foregrounded rather than hidden. It shows that the mechanism is active under stress but not necessarily beneficial. Abrupt betrayal creates a temporal mismatch: the beta channel can sharpen policy selection before the partner model has recalibrated. Lower entropy is then not safer action; it is confident action under a wrong or incomplete model.

This boundary condition is theoretically informative. It predicts that affective deployment should be most helpful when reliability estimates track regime changes at the right timescale, and most harmful when social structure changes faster than the model can update. The gradual-betrayal sensitivity result is consistent with this interpretation: default affect becomes nearly payoff-neutral when the shock is less abrupt.

\subsection{Literature positioning}

The project extends Hesp-style affective precision from a single-agent expected-action-precision account to partner-specific social model-fitness estimates \todocite{Hesp et al. Deeply Felt Affect}. It is complementary to multi-agent active-inference work that explicitly models other agents for strategic planning \todocite{factorised active inference for strategic multi-agent interactions}. Its contribution is not a general multi-agent planning algorithm; it is the partner-local affective precision signal and the behavioral tests that separate model fitness from reward.

The trust task also sits near active-inference models of trust and automation \todocite{active inference trust automation papers}. The manuscript should not claim to be the first active-inference model of trust. The narrower novelty claim is stronger: the model separates partner-specific affective precision from reward and from a dedicated trust scalar, then tests it with model-fitness, deployment-lesion, social-choice, and betrayal-stress readouts.

The H5 variants should be positioned carefully relative to computational psychiatry and active-inference psychiatry reviews \todocite{active inference psychology and psychiatry reviews}. They are not clinical validation. They are parameter perturbations showing that beta dynamics can separate in ways that may be useful for future computational phenotype hypotheses.
